% !TEX program = xelatex
\documentclass{英语试卷}
\begin{document}
\examtitle{南京市 2025 届高三年级上学期学情调研}{英语}
\par 南京市 2025 届高三年级上学期学情调研\par
\par 英语试卷\par
\begin{notice}
\par 1. 答卷前，考生务必将自己的姓名、准考证号填写在答题卡上。\par
\par 2. 回答选择题时，选出每小题答案后，用铅笔把答题卡上对应题目的答案标号涂黑。如需改动，用橡皮擦干净后，再选涂其他答案标号。回答非选择题时，将答案写在答题卡上，写在本试卷上无效。\par
\par 3. 考试结束后，将本试卷和答题卡一并交回。\par
\end{notice}
\parttitle{第一部分 听力（共两节，满分 30 分）}
\par 做题时，先将答案标在试卷上，录音内容结束后，你将有两分钟的时间将试卷上的答案转涂到答题卡上。\par
\sectitle{第一节（共5小题：每小题1.5分，满分7.5分）}
\par 听下面 5 段对话。每段对话后有一个小题，从题中所给的 A、B、C 三个选项中选出最佳选项。听完每段对话后，你都有 10 秒钟的时间来回答有关小题和阅读下一小题。每段对话仅读一遍。\par
\question{例：How much is the shirt?}
\begin{choices}(3)
  \item £ 19.15.
  \item £ 9.18.
  \item £ 9.15. 答案是 C,
\end{choices}
\question{1. What is the man asking about?}
\begin{choices}(3)
  \item Ticket price.
  \item Show duration.
  \item Seat availability.
\end{choices}
\question{2. How does Linda feel toward Tom's proposal?}
\begin{choices}(3)
  \item Grateful.
  \item Guilty.
  \item Worried.
\end{choices}
\question{3. What are the speakers going to do tonight?}
\begin{choices}(3)
  \item Revise a play.
  \item Watch a drama.
  \item Shoot a performance.
\end{choices}
\question{4. What is Jane doing?}
\begin{choices}(3)
  \item Writing a financial report.
  \item Booking a meeting room.
  \item Printing a needed material.
\end{choices}
\question{5. Why will Sarah miss the group presentation?}
\begin{choices}(3)
  \item She got her part covered.
  \item Her flight is delayed.
  \item She didn't plan for it.
\end{choices}
\sectitle{第二节（共15小题：每小题1.5分，满分22.5分）}
\par 听下面 5 段对话或独白。每段对话或独白后有几个小题，从题中所给的 A、B、C 三个选项中选出最佳选项。听每段对话或独白前，你将有时间阅读各个小题，每小题 5 秒钟；听完后，各小题将给出 5 秒钟的作答时间。每段对话或独白读两遍。\par
\par 听第 6 段材料，回答第 6、7 题。\par
\question{6. Where are the speakers?}
\begin{choices}(3)
  \item At a bookstore.
  \item At a library.
  \item At a writing class.
\end{choices}
\question{7. Why does the man refuse to get a half-priced book?}
\begin{choices}(3)
  \item The price is unacceptable.
  \item No other book appeals to him.
  \item He plans to borrow one elsewhere.
\end{choices}
\question{8. What are young passionate architects advised to do?}
\begin{choices}(3)
  \item Wait for the right choice.
  \item Seek professionals' assistance.
  \item Devote themselves to the career.
\end{choices}
\question{9. What is the woman's advice to those uninterested in architecture?}
\begin{choices}(3)
  \item Stick to the path.
  \item Develop their creativity.
  \item Explore a new field.
\end{choices}
\question{10. Who is the man?}
\begin{choices}(3)
  \item A young architect.
  \item A show host.
  \item A career consultant.
\end{choices}
\par 听第 8 段材料，回答第 11 至 13 题，\par
\question{11. What is Tina trying to do?}
\begin{choices}(3)
  \item Find an excuse.
  \item Make a negotiation.
  \item Extend an apology.
\end{choices}
\question{12. Why does Tina want to decline Kelly's party invitation?}
\begin{choices}(3)
  \item She has another engagement.
  \item She is preparing for a test.
  \item She hates loud music.
\end{choices}
\question{13. What policy does Ben follow?}
\begin{choices}(3)
  \item Seize the day.
  \item Be straightforward.
  \item Always attend.
\end{choices}
\par 听第 9 段材料，回答第 14 至 17 题。\par
\question{14. What is the woman doing?}
\begin{choices}(3)
  \item She is slicing meat.
  \item She is rating the service.
  \item She is trying Beijing hot pot.
\end{choices}
\question{15. How should the woman cook the meat?}
\begin{choices}(3)
  \item By holding it with chopsticks.
  \item By boiling it with some Majiang.
  \item By cooking it fully for eight minutes.
\end{choices}
\question{16. Which aspect of Beijing hot pot impresses the woman most?}
\begin{choices}(3)
  \item The soup.
  \item The pot.
  \item The flavor.
\end{choices}
\question{17. What is the probable relationship between the speakers?}
\begin{choices}(3)
  \item Customer and waiter.
  \item Husband and wife.
  \item Tourist and guide.
\end{choices}
\par 听第 10 段材料，回答第 18 至 20 题。\par
\question{18. How many stories will be read today?}
\begin{choices}(3)
  \item 3.
  \item 4.
  \item 5.
\end{choices}
\question{19. What else will the audience get from the channel?}
\begin{choices}(3)
  \item Local news.
  \item Free subscription.
  \item Travel information.
\end{choices}
\question{20. What do we know about the British teenager?}
\begin{choices}(3)
  \item He went abroad alone.
  \item He experimented in Southwest China.
  \item He changed the mind of many villagers.
\end{choices}
\parttitle{第二部分 阅读理解（共两节，满分 50 分）}
\sectitle{第一节（共15小题：每小题2.5分，满分37.5分）}
\par 阅读下列短文，从每题所给的 A、B、C 和 D 四个选项中，选出最佳答案。\par
\passagelabel{A}
\par See the World Together\par
\par Be My Eyes connects blind and low-vision users who want sighted assistance with volunteers and companies willing to help, anywhere in the world, through live video and artificial intelligence.\par
\par Download for IOS\par
\par Blind and low-vision: Access visual information with ease\par
\par Use your smartphone to request video support any time. Just take a picture of a product, and connect with a volunteer or simply use the built-in function Be My AI to describe it.\par
\par Volunteers: Support people when they need it most\par
\par You can support those who are blind or low-vision when they need it most. Join our community and change someone's life today.\par
\par Businesses: Offer accessible customer experiences\par
\par Connecting your customer service with Be My AI can solve up to 90\% of cases automatically and save your time and costs by 60\%. Or route your customers with accessibility needs smoothly through to your specialists.\par
\par Be My Eyes' goal is to make the world more accessible to people with vision loss. We couldn't do this without your help.\par
\par Buy our goods\par
\par Help keep Be My Eyes free by buying our official goods. You'll find T-shirts, stickers and more—and they ship worldwide. Browse\par
\par Be a translator\par
\par Help us make the app accessible for more people in their native language. Translate\par
\par Spread the word\par
\par Invite your friends, family or co-workers. Share 7,536,743 660,439 150+ 180+ Volunteers Blind \& low-vision Countries Languages\par
\question{21. What is the aim of Be My Eyes?}
\begin{choices}(2)
  \item To explore new video and AI applications.
  \item To promote global communication.
  \item To facilitate help for the visually challenged.
  \item To increase the sales of official goods.
\end{choices}
\question{22. Adopting Be My AI, businesses can \blanknum{}.}
\begin{choices}(2)
  \item tailor products to customers' needs
  \item improve efficiency at a lower cost
  \item gain more experienced specialists
  \item conduct automatic team training
\end{choices}
\question{23. What can readers do to support Be My Eyes?}
\begin{choices}(2)
  \item Purchase its service.
  \item Spread native languages.
  \item Upload live video.
  \item Recommend the app.
\end{choices}
\passagelabel{B}
\par In the heart of the Kanien' keha Nation (卡尼恩克哈部族) in Canada, McComber, a native language teacher reflected on a question raised by Maori (毛利) language advocate Sir Karetu: what was his dream for the future of his language?\par
\par Kanien' keha, the Mohawk language, is endangered, and for decades, McComber has worked. hard to operate a native language school. At the turn of the century, as it encountered more barriers, McComber looked around the world for help and inspiration, and Karetu was the answer. "Don't wait for money or approval; just do the work," Karetu told him. It's the same message Karetu carried with him decades ago, when tereo Maori, the Maori language, was also on the line.\par
\par In the early 1970s, young Maori advocates in New Zealand marched to the government, demanding Maori be taught in all public schools. In 1982, Dame Iritana Tawhiwhirangi established Kohanga Reo, an engaging language program where Maori was the only language taught, spoken and used both in the classroom and at home, and the courses focused on their history and culture. "Grandparents became teachers," she says. "The difference-maker was that the Kohanga Reo was entirely community-led. Families raised the money to rent classroom spaces, and volunteers planned and taught classes."\par
\par Translated in English as "language nest", the Kohanga Reo model quickly became an important cultural export and inspired other minority communities worldwide. "Our parents ruled," a student of McComber's nest school says, "They wanted to teach us the beauty of our language, our culture, and who we are. And it has given me the confidence to speak and read with ease in both English and Kanien'keha."\par
\par The model helped communities view the world through their ancestors' eyes, ensuring the survival of their way of life. "It's not just about language," McComber can now answer that question. "It's everything."\par
\question{24. Why did McComber turn to Karetu for help?}
\begin{choices}(2)
  \item Karetu also spoke Kanien' keha.
  \item Karetu once faced a similar crisis.
  \item He needed his financial support.
  \item He desired his tips on language learning.
\end{choices}
\question{25. What contributed to the success of Kohanga Reo?}
\begin{choices}(2)
  \item Expansion of public schools.
  \item Flexibility of the courses.
  \item Rule of the students' parents.
  \item Support from the community.
\end{choices}
\question{26. What does paragraph 4 mainly talk about?}
\begin{choices}(2)
  \item The impact of the Kohanga Reo model.
  \item The importance of parents' involvement.
  \item The procedure for building the nest school.
  \item The case of a student finding his real identity.
\end{choices}
\question{27. By saying "It's everything", McComber suggests \blanknum{}.}
\begin{choices}(2)
  \item guaranteeing the survival of their own schools
  \item keeping their culture for future generations
  \item purifying their unique minority language
  \item practicing the tradition of their ancestors
\end{choices}
\passagelabel{C}
\par Somewhere in the sky, lightning is forming. Although the possibility of getting struck in your lifetime is roughly 1 in 12,000, a human does provide an attractive target for lightning bolts (闪电). Here's what you should expect if you ever find yourself in the path of lightning.\par
\par A lot can happen in the three milliseconds it takes for a lightning bolt to course through your body. Lightning strikes may leave you with deep wounds, often accompanied with third degree burns. Your clothes might even be torn by the explosive force of the surrounding air being superheated to up to 27,700 °C. Blood vessels (血管) exploding from the electric shock and heat often create something called a Lichtenberg figure on your skin. This is a pattern of scars(伤疤) that branches out across your body like a tree.\par
\par But it doesn't end there. You might also be faced with a lifetime of neurological (神经系统的). conditions for reasons that scientists still don't fully understand. It's also possible that you will suffer from pain.\par
\par In some cases, however, a lightning strike can lead to strange super talents. In a blog post, University of Miami neuroscientist Berit Brogaard writes about an incident where a surgeon who was struck by lightning developed an urge to learn to play the piano. He began to compose music he had started hearing in his head since the strike. This type of phenomenon confuses scientists.\par
\par One theory that Brogaard says is that cell death caused by being struck by lightning could cause a one-time flooding of the brain with neurotransmitters(神经递质) that are released from the dying neurons(神经元). This causes a rewiring of neurons, providing access to areas of the brain that were previously inaccessible.\par
\par Cool as it would be, you shouldn't count on lightning turning you into a genius in one flash. Most of its results are painful, and could stay with you for the rest of your life.\par
\question{28. What can we learn from the first two paragraphs?}
\begin{choices}(1)
  \item 1 in 12,000 may die from lightning bolts.
  \item Lightning attacks can heat bodies to 27,700 °
  \item 
  \item Branching scars are typical after lightning attacks.
  \item Humans are the most likely target for lightning bolts.
\end{choices}
\question{29. What happened to the surgeon after the strike?}
\begin{choices}(2)
  \item His unexplored brain areas were activated.
  \item His brain was left in a fixed state of death.
  \item He suffered a lifetime of neurological pain.
  \item He recalled hidden memories about music.
\end{choices}
\question{30. What does the underlined word "This" refer to in paragraph 5?}
\begin{choices}(2)
  \item The cell death due to lightning strikes.
  \item The flooding of brain's neurotransmitters
  \item The rewiring of brain neurons.
  \item The release of dying neurons.
\end{choices}
\question{31. What is the text mainly about?}
\begin{choices}(2)
  \item The process of lightning formation.
  \item The chances of surviving a lightning strike.
  \item The possibility of lightning making a genius.
  \item The consequences of being struck by lightning.
\end{choices}
\passagelabel{D}
\par Can you think of someone who is an excellent communicator? He must have good skills of speaking confidently in various situations, including presentations, group discussions and activities. Also, he should be good at listening to others, understanding group dynamics, and promoting shared decision-making.\par
\par Research on classroom-based group work reveals a paradox(悖论): while cooperative learning can be highly effective, merely placing children in groups often proves unproductive. The solution lies in teaching children how to cooperate effectively. Teachers play a crucial role in this process through their interaction with students.\par
\par To ensure effective group work, teachers must first raise students' awareness of the purpose and value of group activities. More importantly, they need to teach students how to engage in reasoned discussions, known as Exploratory Talk. This involves considering everyone's viewpoints, justifying opinions with reasons, and making collective decisions. In Exploratory Talk, participants engage critically but constructively with each other's ideas, invite contributions, and seek further explanations.\par
\par Successful cooperation requires a reasoned discussion, which can be facilitated by establishing a set. Of Ground Rules' for Exploratory Talk. These rules guide behaviors such as asking for information, seeking and providing reasons, and evaluating proposals. Ground Rules help children work towards joint conclusions, respect opinions, and develop their ideas. For instance, children learn to ask, "What do you think? Why?" and listen attentively to each other's responses, reaching reasoned conclusions.\par
\par Exploratory Talk enhances group activities as a powerful learning tool. Certain activities, like Talking Points, are particularly effective in stimulating such discussions. Talking Points are inspiring statements designed to encourage children to share their knowledge and uncertainties on a specific topic. Research indicates that Talking Points bring about more imaginative and extensive responses than traditional questions. They help children focus on the topic and compare their viewpoints with others'. As they discuss, children realize that their ideas can develop with new evidence and that classmates are valuable resources for new thinking. With proper guidance, the students will soon begin to see their ideas flowing. And gradually they will learn to participate fully not only in education but also in life as well.\par
\question{32. What may help solve the group work paradox?}
\begin{choices}(2)
  \item Active interaction among teachers.
  \item Group learning among students.
  \item Teachers' effective guidance.
  \item Students' deep engagement.
\end{choices}
\question{33. Why should Ground Rules be set for the students?}
\begin{choices}(2)
  \item To popularize group activities.
  \item To build up their individuality.
  \item To stimulate more proposals.
  \item To encourage reasoning in talks.
\end{choices}
\question{34. What can be inferred from the last paragraph?}
\begin{choices}(2)
  \item Students should be challenged with uncertainties.
  \item Traditional questions should be avoided in classes.
  \item Well-designed brainstorming matters in group activities.
  \item Idea exchanges between the students deserve top priority.
\end{choices}
\question{35. Which of the following is a suitable title for the text?}
\begin{choices}(2)
  \item Why putting children in groups doesn't always work
  \item What effects a good communicator has on discussion
  \item Whether teachers should use Exploratory Talk or not
  \item Where students' participation in activities should go
\end{choices}
\sectitle{第二节（共5小题；每小题2.5分，满分12.5分）}
\par 阅读下面短文，从短文后的选项中选出能填入空白处的最佳选项。选项中有两项为多余选项。\par
\par Always the outsider?\par
\par One of the things I used to fear was being the new girl in the office—in the kitchen making coffee while the in-crowd chatted about the weekend. I worried that I wasn't good enough or smart enough. \blanknum{36} I'd struggle for words while trying to look cool. Back at school, I was best friends with Kate and Cathy until we had to walk across the road in twos, and guess who was left out? I've always been the outsider from school days to office life. Or is that the story I have told myself?\par
\par As in so many areas of life, the only thing you can control is your response. I remind myself of what Eleanor Roosevelt said: \blanknum{37} The only person making me fee like an outsider is me. These are the lessons I have learned along the way:\par
\par Starting anywhere new, you just have to find your people, the ones that get you. \blanknum{38} It's OK to take your time and figure out if you're the right fit for each other. Have faith that you will find the right people for you.\par
\par Very often it's our belief systems that tell us we don't fit in. We are replaying the stories from childhood and need to rewrite the script (剧本). \blanknum{39} When are you assuming that you'll be on your own? What fears are you projecting onto yourself?\par
\par Practice self-acceptance. Remind yourself that it's OK to have a different way of operating in the world.\par
\par \blanknum{40} Sometimes, interactions with random strangers can help us feel more confident in all our relationships. For example, volunteering can make us feel connected and give us a sense of belonging.\par
\begin{choices}(1)
  \item Don't charge into friendships, though.
  \item Start noticing your patterns of mentality.
  \item If I was asked a question, panic would set in.
  \item No one can make you feel lower without your agreement.
  \item The feeling of being an outsider would force me to shut up.
  \item Try doing something to take yourself out of your comfort zone.
  \item We all feel a little lost and awkward sometimes—you're not the only one.
\end{choices}
\parttitle{第三部分 语言知识运用（共两节，满分 30 分）}
\sectitle{第一节（共15小题；每小题1分，满分15分）}
\par 阅读下面短文，从短文后各题所给的 A、B、C 和 D 四个选项中，选出可以填入空白处的最佳选项，\par
\par During my holiday in Sri Lanka this year, I experienced a memorable encounter. While \blanknum{41} from Ella to Kandy, we saw salespeople selling goods on the road.\par
\par The \blanknum{42} of this route was the many bends. There were many flower sellers, as each time the road made a bend, I could see them standing, expecting cars to \blanknum{43}.\par
\par But one guy \blanknum{44} ny attention: not the first time, and probably not the second time either, but the third time I noticed something \blanknum{45}. The shirt he was wearing looked like one I had seen before, and at first, I thought this might be the \blanknum{46} dress of the flower sellers. On second thought, I wondered if it could be his \blanknum{47} brother. But when we made another turn, I realized that the same guy was running down, taking shortcuts and \blanknum{48} us all the way!\par
\par It became kind of \blanknum{49}, and at every turn, we were \blanknum{50} whether the guy was still behind us with his flowers. My \blanknum{51} grew: how long would he run after us? After such a long journey, the poor guy started \blanknum{52}.\par
\par I looked at the flowers; they were beautiful and hand-picked. I asked the driver to \blanknum{53}, opened the window, and bought the flowers. The guy was \blanknum{54} and was happy that his \blanknum{55} had paid off!\par
\cz{41}{hiking}{running}{driving}{cycling}
\cz{42}{disadvantage}{characteristic}{function}{destination}
\cz{43}{pull over}{break down}{catch up}{turn around}
\cz{44}{sought}{shifted}{drew}{deserved}
\cz{45}{urgent}{traditional}{beautiful}{familiar}
\cz{46}{fancy}{standard}{casual}{plain}
\cz{47}{twin}{poor}{lost}{industrious}
\cz{48}{observing}{leading}{following}{cheering}
\cz{49}{pleasant}{funny}{inspiring}{meaningful}
\cz{50}{checking}{considering}{doubting}{asking}
\cz{51}{anxiety}{determination}{confidence}{curiosity}
\cz{52}{crying}{sweating}{shouting}{complaining}
\cz{53}{leave}{proceed}{stop}{concentrate}
\cz{54}{smiling}{running}{loading}{bargaining}
\cz{55}{creativity}{service}{dream}{devotion}
\sectitle{第二节（共10小题；每小题1.5分，满分15分）}
\par 阅读下面短文，在空白处填入 1 个适当的单词或括号内单词的正确形式。\par
\par San Francisco, Mar 2, 2024 (SF Standard)— Multiple night markets are set to take place across San Francisco this year, \blanknum{56} (base) on the success of last year's events.\par
\par The Chinatown night market will return first. Starting from March 8, it is scheduled for every second Friday of the month till November. The event organized by Be Chinatown, \blanknum{} nonprofit organization, will offer traditional Chinese foods along with activities.\par
\par Another major night market, the Sunset Night Market with \blanknum{58} (invite) performances, is expected to return in June, according to Joel Engardio, who \blanknum{59} (lead) last year's highly successful event. Engardio said the city is in the process of selecting vendors(摊贩) \blanknum{60} (host) the night market, and he's expecting more night market events throughout the second half of the year, \blanknum{61} (possible) on Irving and Noriega streets. "We want to focus on food," Engardio said. "We want the participants to have the experience just like \blanknum{62} you go to the night market in Bangkok."\par
\par The Indian-themed dance party, the Bhangra and Beats Night Market—also a hit last year with over 10,000 participants—will come back on May 10. The event, part of an effort to breathe life \blanknum{63} the downtown area, will provide a space where family-friendly activities, food and music can be enjoyed.\par
\par "Night markets are an opportunity to show San Francisco to the region and welcome new visitors," the Mayor's Office said in a statement. "These investments are crucial to celebrate the \blanknum{64} (diverse) of the city while contributing to our \blanknum{65} (economy) vitality."\par
\parttitle{第四部分 写作（共两节，满分 40 分）}
\sectitle{第一节（满分15分）}
\par 假定你是李华，上周末你们班在南京进行了为期两天的研学活动。请你给美国朋友 Simon 写一封邮件\par
\par 分享这次经历，内容包括：\par
\par (1) 研学内容；(2) 你的感想。\par
\par 参考词汇: 研学 learning trip\par
\par 注意：(1) 写作词数应为 80 左右；(2) 请按如下格式在答题纸的相应位置作答，\par
\par\vspace{0.4em}
\begin{center}\begin{tabular}{{|l|p{\dimexpr\linewidth-3em-4\tabcolsep}|}}
\hline
Dear Simon, &  \\
Do you have similar experiences? Looking forward to your story. & Yours, Li Hua \\
\hline
\end{tabular}\end{center}
\vspace{0.4em}
\sectitle{第二节（满分25分）}
\par 阅读下面材料，根据其内容和所给段落开头语续写两段，使之构成一篇完整的短文。\par
\par In April last year, sitting in the back of the car, my two daughters, Nia, 10 and Robyn,8, chatted away about the latest electronic devices. "Mum, can I get a new phone for my birthday?" Nia asked. "No, love," I replied. "Can we go to a theme park this year?" Robyn asked. "That's very expensive," I said. Disappointment entered our car like a dark cloud, putting an end to the chat.\par
\par Arriving home, I switched on my laptop, went through pictures taken in Uganda(乌干达). A few months before, I'd been there, working with local schools. As an officer for Teams4u, a charity that provides resources to people in need, I'd seen what life was like there. I went there twice a year and loved it. "You must bring your family next time," one of the local teachers had said to me. My girls had never been out to Africa with me before, but were always wide-eyed as I told my adventures.\par
\par Life in Uganda was very different from what we were used to. Children there had little to eat for lunch. There were no proper desks and chairs in the classroom. I'd been out there to help, but there was still a lot of work to do in some remote schools.\par
\par During dinner, Nia asked me, "Mum, can we dine out?" "Homemade food is boring," Robyn complained. I didn't say a word.\par
\par That night in bed, I thought about the invitation from the teacher in Uganda. By the following morning, I'd made my decision. "Girls, we're going to Uganda! For Two weeks!" I announced. Nia and Robyn were surprised at first but soon became excited! Uganda wasn't a theme park, but it was an adventure anyway.\par
\par So, a few weeks later, in May last year, we landed in the capital city of Kampala. That night, we all slept under mosquito nets (蚊帐) in the blazing (燃烧的) heat—we were worlds away from home. 注意：\par
\par (1) 续写词数应为 150 左右;\par
\par (2) 请按如下格式在答题卡的相应位置作答，\par
\writingpara{Paragraph 1: The next morning, we headed out to a remote school with resources. \blanknum{}}
\writinglines{8}
\writingpara{Paragraph 2: Upon returning to the UK, the change in Nia and Robyn was evident.}
\writinglines{8}
\end{document}
